\clearpage
\section{Full refocusing of the copying machine}
\label{sec:copying-full-refocusing}
The original combined copying machine is now proved refocusing for every
operation in its alphabet. The result holds at every finite depth allowance
and for the unbounded execution. It has no assumed compute-refocusing,
matching-independence or successor-rootedness premise. The search and copying
cases, and reference allocation, were discharged in
Section~\ref{sec:copying-intrinsic-refocusing}. Here we discharge every remaining
case: dereference, assignment, both reference-equality branches, application,
unary negation, addition and subtraction.

\termdef{copying-full-refocusing}
\emph{Full copying refocusing} means that every actual machine transition from
a state rooted by search on its own graph leads to a state rooted by search on
its own new graph. More generally, this theorem holds for any external input
of any fixed interface, not only for closed programs. Every reachable state
therefore has its own pure-search witness. No assumption that execution
eventually terminates is needed.

This closes the coordinate-machine refocusing obligation corresponding to
PDF Lemma~12.4 (HTML Lemma~13,
\href{https://arxiv.org/html/1907.01257v5\#S12.SS1}{Section~12.1}), using the
copying analysis of PDF Lemma~10.4 (HTML Lemma~9). The earlier-search-visit
argument is PDF Lemma~C.5 (HTML Lemma~24); the prefix-replay reasoning follows
the stationary-transition proof of PDF Lemma~C.8 (HTML Lemma~27), in
\href{https://arxiv.org/html/1907.01257v5\#A3}{Appendix~C}.
The general stationary-transition theorem for arbitrary fragments is not
asserted here. Instead, the actual successor descriptions prove exactly the
preservation facts used by each concrete rule.

The dependency map is now: exact search history and absence of earlier visits
to the pending head; a common-occurrence search replay theorem; operation-specific
survivor and incidence facts from the already proved successor certificates;
rootedness for all rewrite families; exhaustive rule recognition; finite and
unbounded refocusing; rootedness throughout executions. Search determinism is
already proved and is used to exclude earlier cycles. No template equivalence
is assumed anywhere in this chain.

Appendix~B.1 removes refocusing from the outstanding machine-law work. The
concrete specimen interface and its locality laws, correspondence with the
paper's explicit focused graphs, the eleven-template argument, semantic
weakening and Figure~22 remain open. Successor determinism remains stated up
to recursive representation; this chapter does not identify named graph
carriers by literal equality. The original terms, Figure~16 translation,
copyability conditions and stable-argument premises are unchanged.
The estimate remains roughly 55--70\% of mathematical effort, with low
confidence: closing refocusing does not measure the size of the remaining
specimen, template and observational proofs.

\Needspace{12\baselineskip}
\subsection{Which edge can a search interaction use?}

\termdef{copying-interactive-primitive}
An \emph{interactive copying primitive} is an instance edge, a term contraction,
or an operation. Case analysis on the actual entering function proves that
whenever it returns a search successor, its primitive belongs to this class.
Atom contractions, weakening, stores, holes and focus labels have no such
entering interaction. Returning interactions already require an operation.

\lean{Copying/PrefixReplay}{8}{9}{Search.Interactive}

\lean{Copying/PrefixReplay}{11}{22}{Search.enter_interactive}

\Needspace{10\baselineskip}
\termdef{copying-used-incidence}
A \emph{used search incidence} records the two ways the token can interact with
an edge: search mode at a source, or backtracking mode at one of its ordered
targets. This is a predicate on the existing graph and token, not an additional
transition rule. The auxiliary position map keeps the mode and applies a
chosen function to the token vertex.

\lean{Copying/PrefixReplay}{25}{27}{Search.Uses}

\lean{Copying/PrefixReplay}{29}{29}{Search.mapPosition}

If a used edge has singleton source $v$, the current history records an
earlier search visit to $v$. In search mode, the token is $v$ itself; in
backtracking mode, the history-stack parent theorem recovers the source's
visit. The same two cases recover the history stack at $v$. Consequently,
if $v$ is an argument of another operation with source $u$, using the first
edge implies an earlier visit to $u$. This last statement will exclude
arithmetic operands and the lambda edge of an application frame from the
prefix before their parent operation.

\lean{Copying/PrefixReplay}{32}{38}{Search.uses_visit}

\lean{Copying/PrefixReplay}{41}{47}{Search.uses_stack}

\lean{Copying/PrefixReplay}{50}{53}{Search.uses_parent_visit}

No interactive primitive is a store: the three cases of the predicate are
disjoint from the store constructor. This excludes the selected store edge
from every interaction which is replayed, even though that edge's target
may change during a read or write.

\lean{Copying/PrefixReplay}{56}{57}{Search.Interactive.ne_store}

\Needspace{12\baselineskip}
\subsection{Replay through a common space of occurrences}

\termdef{copying-common-occurrence-replay}
\emph{Search replay through retained occurrences} compares two graphs using
maps $f:V_G\to W$ and $g:V_H\to W$, where $g$ is injective. The common set
$W$ may include both old host vertices and newly exposed body vertices. There
need not be a total embedding of the source graph into the reduct: arithmetic
and application genuinely delete vertices and edges.

For a search interaction on $G$, suppose its actually used primitive edge
has an edge in $H$ with the same primitive label and with source and target
lists equal after applying $f$ and $g$. If the current tokens agree in $W$,
then the interaction lifts to $H$ and the next tokens agree there too.
For an entering step, injectivity recovers source membership, and naturality
of the entering function supplies the successor. For a return, equality of
mapped singleton sources supplies the corresponding operation source;
equality of the ordered target lists recovers the same argument index.
Naturality of the resume function then supplies the return successor.
Thus the proof preserves the actual search rule and argument order.

\lean{Copying/PrefixReplay}{60}{102}{Search.lift_step}

Consider a prefix ending in search at a pending head $v$, followed by a run
to a rewrite state. At every strict prefix position, the history cannot
already contain a visit to $v$: that would give a positive cycle before a
search normal form, contradicting the exact-length uniqueness proved in
Section~\ref{sec:copying-intrinsic-refocusing}. Assume that every edge used
under this history condition has the retained frame required by the one-step
lemma. Induction lifts the prefix, with exactly its original length. The
induction also extends the source run, so the next application receives the
actual updated history, not an assumed reachability condition.

\lean{Copying/PrefixReplay}{105}{123}{Search.lift_prefix}

\Needspace{10\baselineskip}
To obtain rootedness, split the rooted source request at its earlier search
visit to the request vertex. Apply prefix lifting from the initial input,
whose two images in $W$ agree. The certificate's resumed-token equation and
injectivity of $g$ identify the lifted endpoint with search at the actual
successor token. This theorem leaves explicit frame-survival obligations;
the following sections prove them for the concrete rules.

\lean{Copying/PrefixReplay}{126}{139}{Search.rooted_of_frames}

\Needspace{12\baselineskip}
\subsection{Store writes, reads and identity tests}

The canonical assignment successor exchanges the old and new store-value
attachments, then replaces the assignment head by unit and weakenings.
The pre-head prefix cannot use the assignment edge, since doing so would
record an earlier visit to the pending head. It cannot use the store edge,
which is non-interactive. Every other edge retains its label and source list,
and the previously proved target-exchange theorem preserves its entire ordered
target list. Taking the replacement host map as $f$ and the identity as $g$
therefore proves the canonical write successor rooted. The retained store is
not replaced by a fresh store or by its contents.

\lean{Copying/StoreRefocusing}{9}{27}{Assignment.Match.rooted}

The canonical read successor first grafts the literal sharing construction
at the read head, then exchanges a fresh feedback wire with the old store
contents wire. Every retained edge other than the head and store has unchanged
targets. Indeed, a retained host target cannot equal the fresh graft vertex;
and if it equalled the old contents wire, uniqueness of the incoming edge
would identify its edge with the selected store. The exchange therefore fixes
every entry in the ordered target list. Excluding the head by history and the
store by non-interactivity permits prefix replay and proves the canonical
read successor rooted.

\lean{Copying/StoreRefocusing}{30}{42}{Dereference.Match.retained_target}

\lean{Copying/StoreRefocusing}{45}{62}{Dereference.Match.rooted}

An actual reference comparison supplies its Boolean decision and an
isomorphism from the corresponding canonical successor. That successor is
single-head replacement by the Boolean result and two weakenings, so the
previous head-replacement theorem proves it rooted. Transporting rootedness
along the state isomorphism handles either branch without assuming a decision.

For assignment and dereference, extract the source match from the actual
rewrite, apply the canonical rootedness theorem above, and use the already
proved whole-successor isomorphism. These three public theorems quantify over
arbitrary contexts and actual results. The write's original stable-argument
premise remains in the rewrite relation; no additional premise is imposed on
the caller.

\lean{Copying/StoreRefocusing}{64}{69}{Combined.equality_rooted}

\lean{Copying/StoreRefocusing}{71}{76}{Combined.assignment_rooted}

\lean{Copying/StoreRefocusing}{78}{83}{Combined.dereference_rooted}

\Needspace{12\baselineskip}
\subsection{Application deletes a frame and exposes its body}

An application-opening certificate describes the complete reduct in the
common occurrence space consisting of old host vertices and vertices of the
opened body. Its realization into this space is injective. The pre-head prefix
uses none of the three removed edges. It cannot use the application head
because that would record an earlier visit to the request. It cannot use the
lambda edge because its singleton source is an argument of that application;
the parent-visit lemma would again record the request. It cannot interact with
the box edge, because a recursive box label is not a primitive label.

Every other old edge survives with unchanged ordered sources, targets and
primitive label. The realization's range theorem supplies the surviving edge,
and its incidence equations supply the replay frame. The token equation then
proves rootedness at the exposed body's actual input. This argument retains
the complete recursive body, including nested boxes; it neither executes the
body during opening nor replaces it by an interface summary.

\lean{Copying/ApplicationRefocusing}{8}{36}{Application.OpeningCertificate.rooted}

Every actual application rewrite already yields such a complete opening
certificate. Applying its rootedness theorem discharges application refocusing
in the full host. The original stable argument condition is retained in the
underlying micro-beta rule, even though the prefix-replay argument itself
only needs the certificate's frame and token facts.

\lean{Copying/ApplicationRefocusing}{39}{43}{Combined.application_rooted}

\Needspace{12\baselineskip}
\subsection{Arithmetic deletes its immediate operands}

If a selected argument is a numeric edge, any consumer of its argument wire
is that same edge, by source uniqueness. Thus it has the numeric operation
label. A recognized unary arithmetic request has exactly one target, at index
zero; a binary request has exactly two, at indices zero and one. Inspecting
these indices and applying the numeric-argument fact proves that every edge
adjacent to the arithmetic head is an operation. This is derived from the
actual numeric operands; it is not assumed for arbitrary adjacent edges.

\lean{Copying/ArithmeticRefocusing}{8}{14}{Arithmetic.NumberArgument.consumer}

\lean{Copying/ArithmeticRefocusing}{17}{38}{Arithmetic.Request.deleted_operations}

An arithmetic reduction certificate identifies surviving vertices and edges
with their original occurrences. It retains exactly the edges not adjacent
to the head; the head itself survives with its result label and no targets.
For pre-head replay, the selected head is excluded by the earlier-visit
condition. Every deleted adjacent edge is also excluded: it is an operation
with a singleton source among the head's targets, so interacting with it would
imply an earlier visit to the head. The range theorem therefore supplies a
surviving edge for every interaction that the prefix actually uses.

Use $W=V_G$, $f$ the identity and $g$ the reduction's injective vertex map.
The retained-edge label and incidence equations give the required frames;
the token equation identifies the resumed result. This proves rootedness
even though operand vertices and edges were actually removed.

\lean{Copying/ArithmeticRefocusing}{41}{63}{Arithmetic.Reduction.rooted}

\Needspace{10\baselineskip}
For an arithmetic certificate, uniqueness of the head aligns the original
recognition data with the certificate's selected edge. The request's operand
data then discharges the deleted-operation property, allowing the reduction
rootedness theorem to apply. Finally, every actual arithmetic rewrite supplies
the certificate. This proves the public theorem for unary negation, addition
and subtraction, in arbitrary hosts and without altering their natural-number
semantics.

\lean{Copying/ArithmeticRefocusing}{66}{75}{Arithmetic.Certificate.rooted}

\lean{Copying/ArithmeticRefocusing}{81}{85}{Combined.arithmetic_rooted}

\Needspace{12\baselineskip}
\subsection{The unconditional finite and unbounded theorems}

Every actual contextual rewrite supplies one of the exhaustive original rule
choices. Case analysis applies the proved theorem for allocation, read, write,
comparison, application, copying, unary arithmetic or binary arithmetic.
Consequently every actual rewrite preserves search rootedness. A full-machine
step is either such a rewrite or unchanged-graph search; appending a search
step proves the latter case. This yields preservation from any external input.
Specializing to the sole input of a closed state proves finite refocusing
without a compute-law hypothesis.

\lean{Copying/FullRefocusing}{7}{21}{Combined.rewrite_rooted}

\lean{Copying/FullRefocusing}{24}{28}{Combined.step_rootedFrom}

\lean{Copying/FullRefocusing}{31}{33}{Combined.refocusing}

Induction over a finite run now preserves rootedness from its fixed external
input, using the unconditional one-step theorem at each transition. An
unbounded transition is represented at one finite allowance; applying finite
refocusing proves unbounded refocusing. Another induction gives rootedness
throughout unbounded runs. No common bound on all possible programs is needed,
and no change is made to the existing representation or padding relations.

\lean{Copying/FullRefocusing}{36}{41}{Combined.run_rootedFrom}

\lean{Copying/FullRefocusing}{44}{46}{Combined.unbounded_refocusing}

\lean{Copying/FullRefocusing}{49}{53}{Combined.unbounded_run_rooted}

\termdef{copying-reachable-rootedness}
\emph{Reachable copying rootedness} asserts that every state reached from a
closed program has a pure-search witness on its own current graph. The initial
state has the zero-step witness, and unbounded run preservation proves the
claim for every finite execution length. Given an infinite sequence of actual
steps from an initial program, induction on its index proves the same property
at every position. This last theorem does not assert that an infinite run
exists, or that any particular callback diverges; it shows that refocusing
requires no termination premise.

\lean{Copying/FullRefocusing}{56}{58}{Combined.execution_rooted}

\lean{Copying/FullRefocusing}{61}{67}{Combined.infinite_run_rooted}

\Needspace{12\baselineskip}
\subsection{Nested and store-data regression proofs}

A unary host is the original primitive operation graph with its explicit
single-input, single-output interface. Its entering search step moves the
token to the argument's internal attachment. This is verified directly from
the source and target incidences at either depth form.

\lean{Copying/FullRefocusingExamples}{8}{9}{Combined.FullRefocusingExamples.unaryHost}

\lean{Copying/FullRefocusingExamples}{12}{17}{Combined.FullRefocusingExamples.unary_enters}

Place an application below a reference-allocation host. One host-entry step
and the existing six-step opening execution produce the opened body at an
internal search token. Full run preservation proves this token rooted.
Instantiating with the unchanged nested-lambda translation verifies that
recursive bodies are covered. Instantiating with the original variable-body
translation and its next search step yields a rooted internal copying request;
the variable still triggers actual copying after opening.

\lean{Copying/FullRefocusingExamples}{20}{30}{Combined.FullRefocusingExamples.application_host_rooted}

\lean{Copying/FullRefocusingExamples}{33}{37}{Combined.FullRefocusingExamples.nested_lambda_host_rooted}

\lean{Copying/FullRefocusingExamples}{40}{45}{Combined.FullRefocusingExamples.variable_host_copy_rooted}

Place binary arithmetic below unary negation. The host-entry step followed
by the seven-step binary computation leaves a rooted backtracking token at the
internal argument. The existing finality theorem shows that it is not final:
the enclosing operation must still resume. This checks rootedness across a
real deleting rewrite without confusing a local return with global completion.

\lean{Copying/FullRefocusingExamples}{48}{54}{Combined.FullRefocusingExamples.arithmetic_host_return_rooted}

\lean{Copying/FullRefocusingExamples}{56}{59}{Combined.FullRefocusingExamples.arithmetic_host_return_not_final}

The store regressions combine actual seven-step successful executions with
the new rootedness theorem. A write over an arbitrary unchanged translated old
term has a rooted final result at every finite padding. Both shared-reference
and distinct-reference comparisons also have rooted final results with their
unchanged translated contents, again at every finite padding. The arbitrary
old term may be the existing literal omega translation; these statements do
not claim to prove its divergence. The stored program is not evaluated merely
to establish the write or comparison result.

\lean{Copying/FullRefocusingExamples}{62}{66}{Combined.FullRefocusingExamples.overwrite_rooted_result}

\lean{Copying/FullRefocusingExamples}{69}{80}{Combined.FullRefocusingExamples.equality_rooted_results}

The full build passes 452 jobs. All 2346 axiom audits pass: 6 abstract,
477 graph and 1863 copying declarations, including all 41 declarations in this
chapter. The only axiom dependencies are \texttt{propext},
\texttt{Classical.choice} and \texttt{Quot.sound}. Proposition~12.2 remains
unproved until the specimen, template and final observational obligations in
Appendix~B.1 are completed.
